% Thesis Chapter 8 — Outlook and Conclusion
% .kimi_draft_v3 sidecar
% Zone-scrubbed per R3-1 methodology; all numbers map to 3A/3C.
% No bug-retrospective language.
%
% Seeded from the thesis outline, manuscript results, and post-submission
% experiment roadmap (Routes R-A, R-B, R-C).
%
% Cross-references:
%   Chapter 1 (\ref{chap:hat-instance-overfitting}) — instance overfitting
%   Chapter 2 (\ref{chap:framework}) — simulation framework
%   Chapter 3 (\ref{chap:hat-taxonomy}) — HAT taxonomy
%   Chapter 4 (\ref{chap:failure-mode-atlas}) — failure mode atlas
%   Manuscript: \ref{sec:methodology}, \ref{sec:results}

\chapter{Outlook and Conclusion}
\label{chap:outlook}

\section{Introduction: the thesis as a bridge}

This thesis began with a simple observation.  Organic optoelectronic compute-in-memory devices are invariably introduced through figures of merit---conductance window, retention time, write asymmetry, optical responsivity---that are necessary but not sufficient.  They do not answer the question that ultimately matters for deployment: \emph{what accuracy survives once a modern vision backbone is mapped to an analog array with structured mismatch, finite readout precision, and a sub-linear photoresponse?}  The intervening chapters have built a bridge between the measurement laboratory and the algorithm training loop.  Chapter~\ref{chap:framework} constructed the bridge's truss: a profile-driven simulation framework that turns device characterizations into structured training inputs.  Chapter~\ref{chap:hat-taxonomy} mapped the traffic rules: a taxonomy of hardware-aware training cadences, noise profiles, and transfer guarantees.  Chapter~\ref{chap:failure-mode-atlas} surveyed the load limits: a catalogue of named, reproducible failure modes that bound the deployment envelope.

The preceding chapter constructed the deployment envelope; this chapter steps outside it to propose the next-generation theory and hardware roadmap.

This final chapter has two tasks.  The first is to look forward.  It identifies the experiments, architectures, hardware roadmaps, and theoretical tools that can extend the bridge beyond the scope of the present thesis.  The second is to look backward, integrating the eight chapters into a single narrative arc with three core contributions: (1)~\texttt{compute\-ViT} as a reusable research instrument; (2)~the HAT taxonomy and the discovery of hardware-instance overfitting; and (3)~a quantified deployment envelope for organic CIM.  The outlook is optimistic but grounded; the conclusion is definitive but open-ended---because a thesis, like the bridge it builds, is meant to be crossed.


% =============================================================================
\section{Immediate next steps: thesis-completion experiments}
\label{sec:immediate-next-steps}

Two experiments complete the loops opened in Chapter~\ref{chap:failure-mode-atlas}.  First, \textbf{joint severe-NL and instance-overfitting mitigation}: starting from the canonical Ensemble HAT checkpoint ($86.37\pm1.54\,\%$), apply a weights-only warm-start with MLP blocks retaining $\mathrm{NL}=1.0$ linear surrogates while attention blocks see full $\mathrm{NL}=2.0$ distortion, with D2D resampled every epoch.  The target is $\ge 80\,\%$ fresh-instance accuracy under the $10 \times 5$ protocol; anything below $50\,\%$ would signal a fundamental incompatibility between first-order NL surrogates and distributional training.  Second, \textbf{ImageNet-100 pilot}: test whether the ranking logic (Ensemble HAT $>$ fixed-mask HAT; bounded degradation under correlated D2D) survives dataset scale under the G-AA2 protocol.  Success is qualitative---ranking preservation, not headline numbers.


% =============================================================================
\section{Language-model CIM: transferring HAT principles beyond vision}
\label{sec:language-model-cim}

The dominant transformer workload is the large language model.  Analog CIM offers the same theoretical energy advantage for LLMs, but the mapping is complicated by tiling scale, variable-length attention, and KV-cache storage.  A pilot on a $125$~M-parameter model (Table~\ref{tab:kv-cache-107}) tested whether \emph{selective} terminal-layer analog KV storage could pass a cross-instance robustness gate.

\begin{table}[t]
    \centering
    \caption{\textbf{Fresh-D2D cross-instance robustness of selective analog KV-cache.}  Evaluated on a $125$~M-parameter language model with $5$ independent D2D seeds per evaluation level.  Train D2D $=0.02$ for all selective checkpoints.}
    \label{tab:kv-cache-107}
    \small
    \setlength{\tabcolsep}{5pt}
    \begin{tabular}{@{}llccc@{}}
        \toprule
        \textbf{Checkpoint} & \textbf{Analog layers} & \textbf{Eval D2D} & \textbf{Mean PPL} & \textbf{Std} \\
        \midrule
        Digital baseline        & ---           & ---  & $15.68$  & ---    \\
        \midrule
        last1 selective (train D2D$=0.02$) & $[23]$     & $0.02$ & $18.42$  & $0.02$ \\
        last1 selective (train D2D$=0.02$) & $[23]$     & $0.04$ & $18.55$  & $0.03$ \\
        last1 selective (train D2D$=0.02$) & $[23]$     & $0.05$ & $18.60$  & $0.03$ \\
        \midrule
        last2 selective (train D2D$=0.02$) & $[22,23]$  & $0.02$ & $18.71$  & $0.02$ \\
        last2 selective (train D2D$=0.02$) & $[22,23]$  & $0.04$ & $19.07$  & $0.04$ \\
        last2 selective (train D2D$=0.02$) & $[22,23]$  & $0.05$ & $19.21$  & $0.04$ \\
        \midrule
        all-layer (train D2D$=0.02$)       & all        & $0.05$ & $66.84$  & $5.91$ \\
        \bottomrule
    \end{tabular}
\end{table}

The results yield a clear design rule: \textbf{begin with terminal-layer selective mapping, not all-layer mapping.}  At eval D2D $=0.05$, the all-layer checkpoint catastrophically degrades to $66.84\pm5.91$~PPL ($4.3\times$ the digital baseline), whereas last1 selective degrades only to $18.60\pm0.03$~PPL with cross-instance standard deviation below $0.03$~PPL.  Last1 also outperforms last2 at every tested D2D level, confirming that earlier layers amplify perturbations through multiple downstream paths.  The core thesis principle---train over the deployment distribution, not a single instance---remains valid; the practical challenge is tile-granularity resampling cost, which motivates mismatch resampling only for analog-mapped layers and low-rank surrogate approximations.


% =============================================================================
\section{Hardware-in-the-loop roadmap: from simulation to silicon}
\label{sec:hardware-roadmap}

The ultimate goal is to deploy HAT-trained models on physical organic arrays.  A three-stage roadmap de-risks the transition.  \textbf{Stage~I} replaces first-order surrogates with measured device profiles (spatial covariance, pulse-accurate transients, temperature-dependent retention); the blocker is batch-to-batch variability exceeding within-batch D2D spread.  \textbf{Stage~II} emulates the mixed-signal stack on FPGA to validate control logic; the blocker is resource scaling, since a 256$\times$256 crossbar consumes most of a mid-range FPGA.  \textbf{Stage~III} fabricates a small organic RRAM array (64$\times$64 or 128$\times$128) for functional verification; ecological blockers (humidity, oxygen ingress, contact resistance) dominate.  Calibration infrastructure---per-column ADC calibration, reference-cell tracking, temperature compensation---cuts across all stages and is not yet in the energy model.


% =============================================================================
\section{Theory directions: from empirical recipes to principled guarantees}
\label{sec:theory-directions}

Three directions can transform the empirical recipes into principled guarantees.  First, \textbf{HAT as a regulariser}: epoch-level resampling convolves empirical risk with the mismatch distribution, smoothing sharp minima.  The expected gradient variance bounds the Hessian trace at convergence:
\begin{equation}
    \theta_{\mathrm{ens}}^{\star}
    = \arg\min_{\theta}\;
    \mathbb{E}_{M\sim p(M)}\!\left[
        \frac{1}{N}\sum_{n=1}^{N}
        \ell\!\left(f(x_{n};\theta,M),y_{n}\right)
    \right],
    \label{eq:ens-hat-repeat}
\end{equation}
\begin{equation}
    \mathrm{Tr}\,H(\theta_{\mathrm{ens}}^{\star})
    \;\le\;
    \frac{C}{\sigma_{M}^{2}}\,
    \mathbb{E}\!\left[\|\nabla_{\theta}\ell\|^{2}\right],
    \label{eq:hessian-bound}
\end{equation}
where larger $\sigma_{M}^{2}$ forces a flatter minimum, explaining why Ensemble HAT preserves $86.37\pm1.54\,\%$ fresh-instance accuracy while fixed-mask HAT collapses to $10.00\,\%$.

Second, \textbf{ensemble-frequency effective width}: define $W_{\mathrm{eff}} = E / \tau_{\mathrm{mix}}$, with $E$ training epochs and $\tau_{\mathrm{mix}}$ optimizer mixing time.  Per-epoch resampling gives $W_{\mathrm{eff}}\approx E$; per-batch gives comparable width but shorter mixing, producing the non-stationary gradient field that empirically underperforms.  The inverted-U prediction matches the Chapter~\ref{chap:hat-taxonomy} ablation.
\begin{equation}
    W_{\mathrm{eff}} = \frac{E}{\tau_{\mathrm{mix}}},
    \label{eq:effective-width}
\end{equation}

Third, \textbf{PAC-Bayes bounds for device-instance generalisation}: treating $p(M)$ as a prior and learned weights as posterior $q(\theta)$,
\begin{equation}
    R_{\mathrm{fresh}}(q)
    \;\le\;
    \hat{R}_{\mathrm{ens}}(q)
    \;+\;
    \sqrt{\frac{\mathrm{KL}(q\|p) + \ln\bigl(2\sqrt{E}/\delta\bigr)}{2(E-1)}},
    \label{eq:pac-bayes}
\end{equation}
where the bound tightens as $O(1/\sqrt{E})$, justifying longer training, and the KL term penalises complex posteriors, justifying weight decay.  Future tightening includes Gaussian-process priors for spatial correlation and extension to tiled LLM settings.


% =============================================================================
\section{Conclusion: one story}
\label{sec:conclusion}

A thesis is often read as a sequence of chapters, each with its own figures, equations, and contributions.  But the true product is not the sequence; it is the \emph{story} that runs through it.  This concluding section tells that story in three acts, then gathers its threads into a single claim about what this thesis has built and where it leads.

The story opens with a gap.  Device scientists publish conductance windows and retention curves; system engineers need deployment accuracy.  Chapter~\ref{chap:framework} closed that gap by building \texttt{compute-ViT}, a profile-driven simulation framework in which every physically relevant assumption enters through a structured configuration object.  The framework is not merely an implementation detail; it is a \emph{reusable research instrument}.  A literature-derived prior and an in-house measured profile populate the same fields and drive the same training scripts, so switching from one device generation to the next requires only replacing the profile instance.

With the instrument in hand, Chapter~\ref{chap:hat-instance-overfitting} asked the obvious question: if we train a vision transformer on a noisy analog array, does it survive deployment?  The answer was a stark $10.00\pm0.00\,\%$ collapse.  Standard hardware-aware training, which fixes a single device-to-device mismatch field for the entire training run, overfits that instance as if it were a hidden domain label.  The discovery is not that analog noise hurts accuracy---that was known---but that \emph{robustness to analog noise is not the same as robustness to hardware-instance shift}.  The remedy, Ensemble HAT, is simple in retrospect: resample the mismatch field at epoch boundaries so that the optimizer learns over a \emph{distribution} of arrays rather than memorising one sample from it.  The result---$86.37\pm1.54\,\%$ fresh-instance accuracy---is the empirical anchor of the thesis.  Chapter~\ref{chap:hat-taxonomy} placed that anchor in a broader design space.  Ensemble HAT is not an isolated trick; it is one point in a taxonomy of resampling cadences, noise profiles, and transfer guarantees.  Per-epoch resampling outperforms per-batch because it creates a block-stationary optimisation landscape: stable enough to descend within each block, diverse enough to generalise across blocks.  These distinctions provide a \emph{language} for discussing HAT that did not exist before.

If the first act built the instrument and discovered the central phenomenon, the second act mapped its boundaries.  Chapter~\ref{chap:failure-mode-atlas} organised the empirical landscape as a structured catalogue of named, reproducible breakdown patterns.  The atlas is not a pessimistic document; it is a \emph{deployment envelope}.  It says: if your device operates at $\mathrm{NL}=1.0$ with i.i.d.~mismatch, you can expect $86.37\pm1.54\,\%$ on fresh instances.  At $\rho=0.3$, the mean is $84.57\pm2.39\,\%$; at $\rho=0.5$, it drops modestly to $82.12\pm3.95\,\%$---bounded degradation, not catastrophic collapse.  If retention drift is present but scale recalibration is applied, the accuracy stabilises in the $79.13$--$79.51\,\%$ plateau.  These numbers are contracts between physics and algorithm.  A cross-cutting theme binds the atlas entries together: the primacy of \emph{distributional robustness} over instance-specific rescue.  Standard HAT rescues source-domain accuracy; MLP-linearisation rescues source-domain accuracy under severe NL; but neither survives fresh-instance evaluation.  Only interventions that explicitly average over the deployment distribution produce results that are stable across multiple hardware realisations.  This principle is the methodological legacy of the thesis.

The third act is the outlook developed in the preceding sections of this chapter.  The joint MLP-linear + Ensemble HAT experiment tests whether the two deepest failure modes can be simultaneously repaired.  The ImageNet-100 pilot tests whether the ranking logic survives dataset scale.  The language-model direction tests whether the principles transfer to the dominant AI workload of the coming decade.  The hardware roadmap lays out a staged path from simulation to silicon.  And the theory directions offer a route from empirical recipes to principled guarantees.  These threads are not disjoint futures; they are convergent.  A language-model HAT protocol will need the regularisation theory to choose resampling cadences across thousands of tiles; a chip prototype will need the FPGA stage to validate calibration logic; and every stage will need the profile-driven framework as its common language.  The thesis therefore does not end with a single result; it ends with a \emph{platform}.

It is worth stating the contributions plainly, because a thesis that spans simulation, algorithm, and hardware can easily be misread as three small papers bound together.  It is not.  The contributions are three facets of one story.  First, \texttt{compute-ViT} is a reusable research instrument: the first simulation environment designed specifically for organic optoelectronic CIM that treats device profiles as interchangeable structured inputs, with pluggable physical interfaces, PyTorch autograd integration, and external validation against CrossSim.  Second, the HAT taxonomy formalises the design space of resampling cadences and noise profiles, and the instance-overfitting discovery redefines the evaluation criterion for analog-deployment experiments---the fresh-instance protocol ($10$ arrays $\times$ $5$ Monte Carlo evaluations) is a minimal bar for claiming deployment-grade robustness.  Third, the failure-mode atlas provides the first system-level accuracy contracts for Vision Transformers on organic optoelectronic arrays, telling the materials scientist which device properties matter most, the circuit designer which peripheral overhead is acceptable, and the algorithm designer which training protocol to choose for a given physical regime.

This thesis does not claim that organic optoelectronic compute-in-memory is ready for ImageNet-scale deployment today.  The severe-NL residual gap, the retention plateau, and the environmental sensitivity of organic devices are real blockers that will require years of joint materials--circuit--algorithm research to overcome.  What the thesis claims is something more modest and, perhaps, more durable: that the transition from ``device works'' to ``system works'' requires a new kind of experimental science.  That science must be \emph{profile-driven}, treating device characterisations as structured inputs rather than narrative context.  It must be \emph{distributional}, training over the deployment distribution rather than a single hardware sample.  And it must be \emph{transfer-oriented}, evaluating every intervention on fresh instances before declaring it deployment-grade.  The clean room is quiet and the fabrication tools are precise, but the real world is a distribution of arrays, temperatures, and inference latencies that no single measurement can capture.  Building a bridge between those two worlds is slow, incremental, and often humbling.  Yet every time a trained model survives a fresh hardware instance---every time the accuracy stays above $80\,\%$ on an array it has never seen---the bridge lengthens.  This thesis has lengthened it a little.  The path ahead is to keep building.
